\documentclass[UTF8,11pt]{ctexart}
\usepackage[a4paper,margin=24mm]{geometry}
\usepackage{amsmath,amssymb,amsthm,mathtools,mathrsfs}
\usepackage[all]{xy}
\usepackage{xurl,hyperref,enumitem}
\usepackage{tikz}
\usetikzlibrary{arrows.meta,cd,decorations.markings,calc,patterns}
\hypersetup{hidelinks,breaklinks=true}
\setlength{\emergencystretch}{3em}
\allowdisplaybreaks
\newtheorem*{theorem}{Theorem}
\newtheorem*{thm}{Theorem}
\newtheorem*{lemma}{Lemma}
\newtheorem*{proposition}{Proposition}
\newtheorem*{prop}{Proposition}
\newtheorem*{corollary}{Corollary}
\newtheorem*{cor}{Corollary}
\newtheorem*{claim}{Claim}
\theoremstyle{definition}
\newtheorem*{definition}{Definition}
\newtheorem*{defn}{Definition}
\newtheorem*{remark}{Remark}
\newtheorem*{example}{Example}
\newcommand{\R}{\mathbb R}
\newcommand{\C}{\mathbb C}
\newcommand{\Z}{\mathbb Z}
\newcommand{\Q}{\mathbb Q}
\newcommand{\N}{\mathbb N}
\newcommand{\F}{\mathbb F}
\newcommand{\D}{\mathbb D}
\newcommand{\bB}{\mathbb B}
\newcommand{\bC}{\mathbb C}
\newcommand{\bR}{\mathbb R}
\newcommand{\bZ}{\mathbb Z}
\newcommand{\bN}{\mathbb N}
\newcommand{\bQ}{\mathbb Q}
\newcommand{\bD}{\mathbb D}
\newcommand{\bF}{\mathbb F}
\newcommand{\CP}{\mathbb{CP}}
\newcommand{\RP}{\mathbb{RP}}
\newcommand{\sO}{\mathscr O}
\newcommand{\sI}{\mathscr I}
\newcommand{\sF}{\mathscr F}
\newcommand{\sG}{\mathscr G}
\newcommand{\sH}{\mathscr H}
\newcommand{\sR}{\mathscr R}
\newcommand{\sM}{\mathscr M}
\newcommand{\sV}{\mathscr V}
\newcommand{\sU}{\mathscr U}
\DeclareMathOperator{\Hom}{Hom}
\DeclareMathOperator{\Ext}{Ext}
\DeclareMathOperator{\Tor}{Tor}
\DeclareMathOperator{\Spec}{Spec}
\DeclareMathOperator{\Proj}{Proj}
\DeclareMathOperator{\supp}{supp}
\DeclareMathOperator{\im}{im}
\DeclareMathOperator{\id}{id}
\DeclareMathOperator{\Tr}{Tr}
\DeclareMathOperator{\tr}{tr}
\DeclareMathOperator{\rank}{rank}
\DeclareMathOperator{\ord}{ord}
\DeclareMathOperator{\dist}{dist}
\DeclareMathOperator{\End}{End}
\DeclareMathOperator{\Aut}{Aut}
\DeclareMathOperator{\Gal}{Gal}
\DeclareMathOperator{\vol}{vol}
\DeclareMathOperator{\Cl}{Cl}
\DeclareMathOperator{\coker}{coker}
\DeclareMathOperator{\codim}{codim}
\DeclareMathOperator{\divergence}{div}
\DeclareMathOperator{\Ric}{Ric}
\DeclareMathOperator{\grad}{grad}
\DeclareMathOperator{\Hess}{Hess}
\newcommand{\abs}[1]{\left\lvert#1\right\rvert}
\newcommand{\sabs}[1]{\lvert#1\rvert}
\newcommand{\babs}[1]{\bigl\lvert#1\bigr\rvert}
\newcommand{\Babs}[1]{\Bigl\lvert#1\Bigr\rvert}
\newcommand{\bbabs}[1]{\biggl\lvert#1\biggr\rvert}
\newcommand{\BBabs}[1]{\Biggl\lvert#1\Biggr\rvert}
\newcommand{\norm}[1]{\left\lVert#1\right\rVert}
\newcommand{\snorm}[1]{\lVert#1\rVert}
\newcommand{\bnorm}[1]{\bigl\lVert#1\bigr\rVert}
\newcommand{\Bnorm}[1]{\Bigl\lVert#1\Bigr\rVert}
\newcommand{\bbnorm}[1]{\biggl\lVert#1\biggr\rVert}
\newcommand{\BBnorm}[1]{\Biggl\lVert#1\Biggr\rVert}
\newcommand{\widebar}[1]{\overline{#1}}
\newcommand{\nicefrac}[2]{\frac{#1}{#2}}
\newcommand{\linnprod}[2]{\langle#1,#2\rangle}
\newcommand{\myindex}[1]{#1}
\newcommand{\glsadd}[1]{}
\newcommand{\avoidbreak}{}
\newcommand{\sourceref}[1]{\textnormal{[原文：\nolinkurl{#1}]}}
\newcommand{\thmref}[1]{\sourceref{#1}}
\newcommand{\lemmaref}[1]{\sourceref{#1}}
\newcommand{\propref}[1]{\sourceref{#1}}
\newcommand{\corref}[1]{\sourceref{#1}}
\newcommand{\defnref}[1]{\sourceref{#1}}
\newcommand{\exampleref}[1]{\sourceref{#1}}
\newcommand{\exerciseref}[1]{\sourceref{#1}}
\newcommand{\figureref}[1]{\sourceref{#1}}
\newcommand{\sectionref}[1]{\sourceref{#1}}
\newcommand{\appendixref}[1]{\sourceref{#1}}
\newcommand{\stacksref}[2]{\href{https://stacks.math.columbia.edu/tag/#1}{#2 [#1]}}


\newcommand{\E}{\mathbb E}
\newcommand{\Prob}{\mathbb P}
\newcommand{\Reff}{\mathcal R_{\mathrm{eff}}}
\newcommand{\eps}{\varepsilon}
\DeclareMathOperator{\diam}{diam}
\DeclareMathOperator{\Var}{Var}
\DeclareMathOperator{\Crit}{Crit}
\newcommand{\dd}{\,\mathrm d}

\begin{document}
\section*{069\quad 完美域上有限生成扩张的可分超越基}
\subsection*{来源与计数目标}
J. S. Milne, \emph{Fields and Galois Theory}, version 5.10, September 2022，Theorem 9.27，印刷p.120；同时收录其调用的 Theorem 5.1 原始元论证，pp.61--62。访问日期2026-09-28。
\par\url{https://www.jmilne.org/math/CourseNotes/FT.pdf}
\par\textbf{本文件只计一个问题：}每个完美域的有限生成扩张都存在一个超越基，使得相应代数扩张有限可分。
\subsection*{作者命题与人类证明}
\begin{theorem}[5.1: primitive element theorem]
Let $E=F[\alpha_1,\ldots,\alpha_r]$ be a finite extension of $F$, and assume that $\alpha_2,\ldots,\alpha_r$ are separable over $F$ (but not necessarily $\alpha_1$). Then there exists a $\gamma\in E$ such that $E=F[\gamma]$.
\end{theorem}
\begin{proof}
For finite fields, we proved this in 4.19. Hence we may assume $F$ to be infinite. It suffices to prove the statement for $r=2$, for then
\[F[\alpha_1,\alpha_2,\ldots,\alpha_r]=F[\alpha'_1,\alpha_3,\ldots,\alpha_r]=F[\alpha''_1,\alpha_4,\ldots,\alpha_r]=\cdots.\]
Thus let $E=F[\alpha,\beta]$ with $\beta$ separable over $F$. Let $f$ and $g$ be the minimal polynomials of $\alpha$ and $\beta$ over $F$, and let $L$ be a splitting field for $fg$ containing $E$. Let $\alpha_1=\alpha,\ldots,\alpha_s$ be the roots of $f$ in $L$, and let $\beta_1=\beta,\beta_2,\ldots,\beta_t$ be the roots of $g$. For $j\ne1$, $\beta_j\ne\beta$, and so the the equation
\[\alpha_i+X\beta_j=\alpha+X\beta\]
has exactly one solution, namely, $X=(\alpha_i-\alpha)/(\beta-\beta_j)$. If we choose a $c\in F$ different from any of these solutions (using that $F$ is infinite), then
\[\alpha_i+c\beta_j\ne\alpha+c\beta\quad\text{unless }i=1=j.\]
Let $\gamma=\alpha+c\beta$. I claim that $F[\alpha,\beta]=F[\gamma]$.

The polynomials $g(X)$ and $f(\gamma-cX)$ have coefficients in $F[\gamma]$, and have $\beta$ as a root:
\[g(\beta)=0,\qquad f(\gamma-c\beta)=f(\alpha)=0.\]
In fact, $\beta$ is their only common root, because we chose $c$ so that $\gamma-c\beta_j\ne\alpha_i$ unless $i=1=j$. Therefore
\[\gcd(g(X),f(\gamma-cX))=X-\beta.\]
Here we computed the gcd in $L[X]$, but this is equal to the gcd computed in $F[\gamma][X]$ (Proposition 2.17). Hence $\beta\in F[\gamma]$, and this implies that $\alpha=\gamma-c\beta$ also lies in $F[\gamma]$. This proves the claim.
\end{proof}

\paragraph{Separating transcendence bases.}
Let $E\supset F$ be fields with $E$ finitely generated over $F$. A subset $\{x_1,\ldots,x_d\}$ of $E$ is a \emph{separating transcendence basis} for $E/F$ if it is algebraically independent over $F$ and $E$ is a finite separable extension of $F(x_1,\ldots,x_d)$.

\begin{theorem}[9.27]
If $F$ is perfect, then every finitely generated extension $E$ of $F$ admits a separating transcendence basis over $F$.
\end{theorem}
\begin{proof}
If $F$ has characteristic zero, then every transcendence basis is separating, and so the statement becomes that of Theorem 9.10. Thus, we may assume $F$ has characteristic $p\ne0$. Because $F$ is perfect, every polynomial in $X_1^p,\ldots,X_n^p$ with coefficients in $F$ is a $p$th power in $F[X_1,\ldots,X_n]$:
\[\sum a_{i_1\cdots i_n}X_1^{i_1p}\cdots X_n^{i_np}
=\left(\sum a_{i_1\cdots i_n}^{1/p}X_1^{i_1}\cdots X_n^{i_n}\right)^p.\]
Let $E=F(x_1,\ldots,x_n)$, and assume that $n>d+1$, where $d$ is the transcendence degree of $E$ over $F$. After renumbering, we may suppose that $x_1,\ldots,x_d$ are algebraically independent (9.9). Then $f(x_1,\ldots,x_{d+1})=0$ for some nonzero irreducible polynomial $f(X_1,\ldots,X_{d+1})$ with coefficients in $F$. Not all $\partial f/\partial X_i$ are zero, for otherwise $f$ would be a polynomial in $X_1^p,\ldots,X_{d+1}^p$, which implies that it is a $p$th power. After renumbering $x_1,\ldots,x_{d+1}$, we may suppose that $\partial f/\partial X_{d+1}\ne0$.

Then $x_{d+1}$ is separably algebraic over $F(x_1,\ldots,x_d)$ and $F(x_1,\ldots,x_{d+1},x_{d+2})$ is algebraic over $F(x_1,\ldots,x_{d+1})$, hence over $F(x_1,\ldots,x_d)$ (1.31), and so, by the primitive element theorem (5.1), there is an element $y$ such that
\[F(x_1,\ldots,x_{d+2})=F(x_1,\ldots,x_d,y).\]
Thus $E$ is generated by $n-1$ elements (as a field containing $F$). After repeating the process, possibly several times, we will have $E=F(z_1,\ldots,z_{d+1})$ with $z_{d+1}$ separable over $F(z_1,\ldots,z_d)$.
\end{proof}
\subsection*{整理说明（非作者正文）}
本条只计9.27；5.1是支撑，不另计为“原始元定理”新题。原语言、两个证明的内部顺序与重复单词照录。未收入其后9.28的几何旁注，不把旁注中的额外措辞并入本题。

标准先修为超越基抽取与有限生成扩张的超越次数、不可约多项式的导数可分性判据、代数扩张的传递性、域扩张不改变多项式最大公因式。5.1有限域的情形引用有限域乘法群循环性；本题实质构造采用5.1中的无限域论证。

难度依据：完美性并不使任意超越基自动可分；作者要利用Frobenius与形式偏导识别可以调整的生成元，继而以原始元构造逐步减少代数生成元，得到所需超越基。正文来自PDF可提取文本的TeX转录，没有另造微分模证明，也未声称取得PDF原件。
\subsection*{可修改的简短标注（非作者正文）}
$c$：函数域扩张的可分超越基存在性。

$a$：完美域；Frobenius；不可约多项式；形式偏导；可分性；原始元；生成元消去；超越次数。

\texttt{cross\_theory\_status = yes}。从完美域上的幂结构推出某个形式偏导非零，转为一个生成元的可分性，再通过原始元的线性组合构造调整整个域的生成表示。位置：9.27的正特征论证与5.1。
标签为非穷尽、可修改初稿。
\subsection*{许可与修改记录}
原作者及作品见来源栏；来源采用 CC BY-NC-SA 4.0。本题证摘录和附加整理沿用该许可。2026-09-28：增加题号、来源定位、独立排版、整理说明和初稿标注；数学内容的取材范围及排版变化见整理说明。
\end{document}
